\begin{table}[htbp]
\centering
\footnotesize
\setlength{\tabcolsep}{3.5pt}
\caption{The Distributive Trajectory, Capacity Utilization, Profit Rate, and Net Capital Accumulation (1960--1975)}
\label{tab:profit_rate_accumulation_1960_1975}
\resizebox{\textwidth}{!}{%
\begin{tabular}{lcccccccl}
\toprule
\textbf{Year} & $\bm{\omega}$ \textbf{(\%)} & $\bm{u}$ \textbf{(\%)} & $\bm{\mu}$ \textbf{(\%)} & $\bm{Y/K}$ & $\bm{r}$ \textbf{(\%)} & $\bm{\chi}$ & $\bm{g^n_K}$ \textbf{(\%)} & \textbf{Regime Dynamic} \\
\midrule
1960 & 60.9\% & 6.6\% & 86.9\% & 0.2580 & 10.08\% & 0.421 & 4.44\% & Pre-reform stabilization baseline \\
1961 & 65.9\% & 7.5\% & 87.3\% & 0.2424 & 8.27\% & 0.534 & 4.62\% & Early ISI expansion \\
1962 & 59.4\% & 7.4\% & 87.9\% & 0.2593 & 10.52\% & 0.421 & 4.63\% & Exchange rate devaluation \\
1963 & 55.8\% & 7.0\% & 89.8\% & 0.2670 & 11.81\% & 0.423 & 5.25\% & Pre-electoral demand expansion \\
1964 & 67.7\% & 6.5\% & 88.4\% & 0.2525 & 8.16\% & 0.563 & 4.82\% & Frei Montalva election; wage surge \\
1965 & 63.8\% & 5.9\% & 86.7\% & 0.2374 & 8.59\% & 0.420 & 3.75\% & Initial anti-inflationary program \\
1966 & 55.3\% & 5.6\% & 93.4\% & 0.2450 & 10.94\% & 0.323 & 3.66\% & Output peak under Frei \\
1967 & 57.9\% & 4.2\% & 93.5\% & 0.2410 & 10.14\% & 0.356 & 3.75\% & Agrarian reform promulgated; profit squeeze \\
1968 & 53.9\% & 4.4\% & 93.7\% & 0.2452 & 11.29\% & 0.337 & 3.96\% & Drought and fiscal austerity \\
1969 & 50.1\% & 5.0\% & 94.1\% & 0.2560 & 12.79\% & 0.288 & 3.82\% & Pre-electoral stagflation; \textit{El Tacnazo} \\
1970 & \textbf{50.9\%} & 5.9\% & 92.8\% & 0.2491 & \textbf{12.23\%} & 0.332 & \textbf{4.23\%} & Allende election; inherited slack baseline \\
1971 & \textbf{56.2\%} & 5.2\% & \textbf{98.4\%} & 0.2427 & \textbf{10.62\%} & 0.304 & \textbf{3.34\%} & Idle capacity mobilized; wage expansion \\
1972 & \textbf{69.0\%} & \textbf{4.0\%} & \textbf{96.0\%} & 0.2121 & \textbf{6.57\%} & \textbf{0.277} & \textbf{1.85\%} & Profit squeeze peak; Bosses' Lockout \\
1973 & \textbf{54.5\%} & 4.8\% & 89.5\% & 0.2004 & \textbf{9.12\%} & \textbf{0.184} & \textbf{1.71\%} & Coup d'état; DL 522 price deregulation \\
1974 & \textbf{39.4\%} & 9.1\% & 88.8\% & 0.1762 & \textbf{10.67\%} & 0.218 & \textbf{2.38\%} & Junta wage compression; union ban \\
1975 & 49.3\% & \textbf{18.0\%} & \textbf{76.6\%} & 0.1289 & 6.54\% & 0.183 & \textbf{1.21\%} & Chicago shock therapy depression \\
\bottomrule
\end{tabular}%
}
\begin{minipage}{0.96\textwidth}
\vspace{0.3em}
\footnotesize
\textit{Source:} Author's calculations based on \citet{Astorga2023}, Clio-Lab PUC / \citet{DiazLudersWagner2016}, and Chapter~2 dataset. Following Chapter~2, productive net capital stock ($K_t \equiv K_{\text{ME}, t} + K_{\text{NRC}, t}$) strictly excludes Residential Construction ($RC$). Structural capacity utilization ($\mu \equiv Y/Y^p$) is estimated via the CMPR IM-OLS cointegrating framework of Chapter~2 (\texttt{Capacity-Utilization-US\_Chile}); $u$ is the open unemployment rate. Implicit annual depreciation is $D_t \equiv K_{t-1} + I_t - K_t$ with depreciation rate $\delta_t \equiv D_t / K_{t-1}$. The aggregate profit rate satisfies the Weisskopf-Marxian decomposition $r_t \equiv (1 - \omega_t) Y_t / (P_{K, t} K_t) = (1 - \omega_t) \mu_t \sigma_t$, where $\sigma_t \equiv Y^p_t / K_t$ is potential capital productivity. Net capital accumulation satisfies the classical Cambridge-Marxian identity $g^n_{K, t} \equiv (I_t - D_t) / K_{t-1} = \chi_t \cdot r_t$, where $\chi_t \equiv I^{\text{net}}_t / \Pi_t$ is the capitalist re-capitalization propensity out of surplus value $\Pi_t \equiv (1 - \omega_t) Y_t$.
\end{minipage}
\end{table}
